\section{Decomposing the actual classical and ordinary spaces}
\label{sec:classical-parts-v43}

This chapter proves the original \texttt{TClParts.clParts\_alt} and its
version for an arbitrary predicate on the parts.  The input is the original
\texttt{AutomorphicHeckeSet}: vectors in the big automorphic space, with
some finite level and some Iwahori depth, are algebraic sums of the stated
parts.  The conclusion preserves the \emph{given} level, finite-part
character, depth, and ordinary generalized-eigenvalue condition.  In
particular, no passage to a closure or to a different automorphic space is
made.  The proof of the big automorphic decomposition, and the arithmetic
construction of the Galois data decorating each part, are distinct inputs.

All source correspondences use commit
\texttt{7cbf7a0c32fca1badf05931879a37c3f8a2b03b4}.
The notation is that of FNL, OST, and ASP.  Fix the original datum
\(H\), satisfying \(H.\mathrm{Ax}\), a complete Noetherian local
coefficient ring \(A\), its original weight character \(\theta\),
an arithmetic point \(x:A\to\overline{\mathbf Q}_\ell\),
the fixed field isomorphism \(\iota\), an admitted weight \(\lambda\),
and the original algebraic representation \(\xi\) on \(\mathbf C^d\).
Write
\[
 \mathcal M=\{F:G^a\times G^\ell\longrightarrow
                  \mathbf C^{m}\otimes\mathbf C^d\}.
\]
This is the same ambient function space as the source's iterated function
type.  Let \(\rho_a\) denote the away action with coefficient factor
\(\tau_C\), and \(\rho_p\) the product local right action.  For an
auxiliary level \(K=(K_q)\), depth \(e\), and the original finite-part
character \(\phi=\mathrm{finPart}(\theta,x,\lambda)\), put
\(C_{K,e}=H.\mathrm{ClAt}(\theta,\iota,x,\lambda,\xi,K,e)\).
The ordinary space is denoted \(O_{K,e}\).

An original \emph{part} \(Y\) is the actual subspace in a
\texttt{PartAltAlt} record.  Among its fields are \(Y\subset H.\mathrm{Big}\),
nonvanishing, and local isotypic presentations at every auxiliary place
and every place above \(\ell\).  Its other fields retain the spherical
eigenvalues, Galois representation, and local comparison conditions.  This
chapter never modifies any of these fields.  A family may be selected by
\(P(\lambda,r_Y)\), as in \texttt{clParts\_alt}, or by a predicate on the
whole record, as needed in \texttt{OrdPartPt}.

\subsection{Averaging while preserving all scalar equations}

% BEGIN CFD1
\paragraph{CFD1. The scalar profile of one vector.}
\label{claim:CFD1}
Let \(\rho:G\to\operatorname{GL}(M)\) be a representation over
\(\mathbf C\), \(K\le G\), and \(F\in M\).  Define
\[
 E_K(F)=\{y\in M: \rho(k)F=aF\Longrightarrow \rho(k)y=ay
                       \text{ for every }k\in K,\ a\in\mathbf C\}.
\]
This is a linear subspace.  Every endomorphism commuting with all
\(\rho(k)\), \(k\in K\), preserves it.  If \(F\ne0\), then
\[
 S=\{k\in K:\rho(k)F\in\mathbf C F\}
\]
is a subgroup and there is a unique character
\(\chi:S\to\mathbf C^\times\) with \(\rho(s)F=\chi(s)F\).
Moreover
\[
 E_K(F)=\{y:\rho(s)y=\chi(s)y\ (s\in S)\}.
\]

For the first assertions, apply linearity to each displayed scalar
equation, and move a commuting endomorphism across \(\rho(k)\).
When \(F\ne0\), the scalar in such an equation is unique and nonzero,
since \(\rho(k)\) is invertible.  Applying \(\rho(s)\rho(t)\) to
\(F\) proves multiplicativity of the scalar; applying \(\rho(s^{-1})\)
proves the inverse assertion.  These calculations prove both the subgroup
claim and the description of \(E_K(F)\).  The nonzero hypothesis is used
only for this description; the definition itself also covers \(F=0\).
% END CFD1

% BEGIN CFD2
\paragraph{CFD2. Finite-index averaging without a normality assumption.}
\label{claim:CFD2}
Let \(\pi:S\to\operatorname{GL}(M)\), and let \(U\le S\) have
finite index.  Choose representatives \(r_q\) of the left cosets
\(S/U\), and define the endomorphism
\[
 P=\frac1{[S:U]}\sum_{q\in S/U}\pi(r_q).
\]
If \(y\) is \(U\)-fixed, then \(Py\) is \(S\)-fixed and does not
depend on the chosen representatives.  If \(y\) is already \(S\)-fixed,
then \(Py=y\).  Finally \(P\) preserves every \(S\)-stable subspace.

Replacing \(r_q\) by \(r_qu\), \(u\in U\), does not change its
action on a \(U\)-fixed vector.  For \(s\in S\), the representatives
\(sr_q\) enumerate the same left cosets, so \(\pi(s)Py=Py\).
The other assertions follow term by term.  No assertion that \(P\)
projects all of \(M\) is needed: invariance of \(Py\) is asserted on
the \(U\)-fixed domain.  In particular the argument does not incorrectly
treat a nonnormal \(U\) as a quotient group.
% END CFD2

% BEGIN CFD3
\paragraph{CFD3. Refinement of a finite sum to the scalar profile.}
\label{claim:CFD3}
Let \((Y_i)_{i\in I}\) be subspaces of \(M\), stable under \(K\).
Suppose every \(y\in Y_i\) is fixed by a subgroup \(U_y\le G\) for
which \([K:K\cap U_y]<\infty\).  If
\(F\in\sum_iY_i\), then
\[
                         F\in\sum_i(Y_i\cap E_K(F)).
\]
The sum here is algebraic, so each membership involves only finitely many
summands.

There is nothing to prove for \(F=0\).  Otherwise write
\(F=\sum_{i\in I_0}y_i\), with \(I_0\) finite.  Let \(S,\chi\)
be as in CFD1, and put \(U=S\cap\bigcap_{i\in I_0}U_{y_i}\).
Its index in \(S\) is finite: restriction embeds each coset set into the
corresponding finite coset set for \(K\), and intersections inject into
products of these coset sets.  Every \(u\in U\) fixes all \(y_i\), hence
fixes their nonzero sum, so \(\chi(u)=1\).  Apply CFD2 to the twisted
representation \(\pi(s)=\chi(s)^{-1}\rho(s)\) and this \(U\).
Each \(y_i\) is \(\pi(U)\)-fixed; therefore \(Py_i\) belongs to
\(Y_i\cap E_K(F)\).  Since \(F\) is \(\pi(S)\)-fixed,
\(F=PF=\sum_iPy_i\).  This proves the assertion without assuming that
the original sum of the \(Y_i\)'s is direct.
% END CFD3

% BEGIN CFD4
\paragraph{CFD4. Simultaneous refinement for commuting local factors.}
\label{claim:CFD4}
Suppose a finite family of groups \(G_j\) acts on \(M\), the actions
of distinct factors commute, and every \(Y_i\) satisfies the hypotheses
of CFD3 for each \(K_j\le G_j\).  Then
\[
 F\in\sum_iY_i\quad\Longrightarrow\quad
 F\in\sum_i\left(Y_i\cap\bigcap_j E_{K_j}(F)\right).
\]

Apply CFD3 to the first factor.  By CFD1 and commutation, every later
factor preserves the scalar-profile subspace already imposed.  Its
intersection with a \(Y_i\) is still stable, and its vectors still have
the required finite-index stabilizers, since they are vectors of \(Y_i\).
Induction on the finite set of factors proves the displayed formula.
An empty family gives the initial algebraic sum itself.
% END CFD4

% BEGIN CFD5
\paragraph{CFD5. Local isotypic sums supply the required stabilizers.}
\label{claim:CFD5}
Let \(G\) be a topological group, \(K\le G\) compact, and let
\(Y\) be an algebraic sum of images of equivariant maps from one smooth
representation \(\sigma\).  Then \(Y\) is \(G\)-stable and every
\(y\in Y\) is fixed by a subgroup \(U\) with
\([K:K\cap U]<\infty\).

Write \(y\) as a finite sum of images of vectors of \(\sigma\).
Smoothness gives an open stabilizer for each vector; their finite
intersection \(U\) fixes \(y\).  The open cosets of \(K\cap U\)
cover compact \(K\), so finitely many suffice.  Equivariance proves
stability term by term.  This is exactly the content of the original
\texttt{IsoSum} fields of a part at each auxiliary place and each place
above \(\ell\).  Irreducibility of \(\sigma\) is not required for
this particular implication, although it remains part of the original
part record.
% END CFD5

\subsection{The original level and ordinary operators}

% BEGIN CFD6
\paragraph{CFD6. Commutation on the original ambient function space.}
\label{claim:CFD6}
The actions at different auxiliary places commute, the actions at
different places above \(\ell\) commute, and every auxiliary action
commutes with every action above \(\ell\).  Every original part is
stable under the whole product \(G^\ell\).

For the first assertion, the \(Q\)-\(Q\) field of \(H.\mathrm{Ax}\)
identifies the original sections as commuting elements in \(U_b\).
Apply the coefficient action \(\rho_a\) to that identity.  For the
second assertion, the local elements lie in distinct coordinates of
\(G^\ell\), hence commute.  For the mixed assertion, on each coordinate
of a function the two expressions are
\[
  \tau_C(u)F\bigl(g(u,1)(1,b)\bigr),\qquad
  \tau_C(u)F\bigl(g(1,b)(u,1)\bigr),
\]
which agree.  These are the original actions; no abstract conjugacy is
used.  FNL12 verifies the relevant section identities for the actual
unitary datum.  Local stability of a part follows from CFD5, and finite
coordinate factorization gives stability under all of \(G^\ell\).
% END CFD6

% BEGIN CFD7
\paragraph{CFD7. Exact classical equations imply the required Iwahori fixedness.}
\label{claim:CFD7}
The original classical space is exactly the space of \(F\in H.\mathrm{Big}\)
satisfying
\begin{align*}
 \rho_a(u)F&=F &&(u\in H.\mathrm{lev}(K)),\\
 \rho_p(b)F&=F &&(b\in N\text{ or }b\in N^-(e)),\\
 \rho_p(e_V(v,i,y))F&=\iota(\phi(v,i,y))^{-1}F
                  &&(y\in\mathcal O_{L_{w(v)}}^\times).
\end{align*}
Suppose the specialization datum has depth \(c_0\), so that
\(\phi(v,i,y)=1\) on the original \(U_c(c_0+1)\).  For every
\(e\ge c_0\), a vector of \(C_{K,e}\) is fixed by \(I_1(e)\).

The first assertion is expansion of the original definition, as already
identified in ASP5.  The subgroup \(I_1(e)\) is generated by
\(N,T_c(e),N^-(e)\).  The first and last are fixed by the displayed
equations.  Each torus coordinate of \(T_c(e)\) is a product of the
original elementary diagonal elements \(e_V(v,i,y_i)\), with
\(y_i\in U_c(e+1)\subset U_c(c_0+1)\).  The finite-part character
is therefore one on each factor.  This fixes the middle generator too.
The factorization and the shift by one are the exact GL and product
Iwahori identities in OST26--OST27 and ASP4, not a replacement convention
for depth.
% END CFD7

% BEGIN CFD8
\paragraph{CFD8. The actual distinguished operator preserves all required spaces.}
\label{claim:CFD8}
Let \(T=H.\mathrm{uzC}(\iota,\lambda,d)\), with its original
\(\mathrm{lowVal}\) normalization.  If \(e\ge c_0\), then \(T\)
preserves \(C_{K,e}\).  It also preserves every original part \(Y\).

On the product action, \(T\) is its stated nonzero scalar times the
finite \(N\)-coset sum defining \(U_z\).  CFD6 shows that each part is
stable under every term.  The big space and its away level condition are
also preserved term by term: finite right translation preserves rational
equivariance, and the away action commutes with it.

For the remaining classical conditions, CFD7 gives \(I_1(e)\)-fixedness.
The original Iwahori level calculation, OST5--OST6, shows that \(U_z\)
preserves this fixedness at depth \(e\).  The same double-coset calculation
commutes \(U_z\) with each elementary integral torus element on
\(N\)-fixed vectors.  Thus its prescribed scalar equation is preserved.
Together these give every equation in CFD7.  Multiplication by the
original normalization scalar does not change stability.  This also
agrees with the positive-operator identification in ASP6.
% END CFD8

% BEGIN CFD9
\paragraph{CFD9. A decomposition in big parts preserves the given classical space.}
\label{claim:CFD9}
Let \((Y_i)\) be any family of original parts and let
\(F\in C_{K,e}\cap\sum_iY_i\).  Then
\[
                         F\in\sum_i(Y_i\cap C_{K,e}).
\]
This implication does not require \(e\ge c_0\).

Apply CFD4 first to the auxiliary groups with their compact subgroups
\(K_q\), and then to the factors above \(\ell\) with
\(\mathrm{GL}_n(\mathcal O_{L_{w(v)}})\).  CFD5 gives finite-index
stabilizers and CFD6 gives all commutations.  We obtain a finite sum
\(F=\sum_iF_i\), where \(F_i\in Y_i\) and every \(F_i\) satisfies
every scalar equation of \(F\) under those compact subgroups.

Since \(Y_i\subset H.\mathrm{Big}\), the base away subgroup is already
fixed on \(F_i\).  The extra \(K_q\)-conditions are inherited from
the scalar profile, with scalar one.  They generate
\(H.\mathrm{lev}(K)\) together with that base subgroup.  The original
\(N\), \(N^-(e)\), and integral torus all lie in the above maximal
compact product, by the GL Iwahori identities of OST26.  Their fixedness
or prescribed finite-part scalar equations are inherited from \(F\).
Thus all equations of CFD7 hold for each \(F_i\), proving the claim.
Neither the level \(K\) nor the depth \(e\) has changed.
% END CFD9

% BEGIN CFD10
\paragraph{CFD10. Polynomial projection preserves a decomposition into invariant pieces.}
\label{claim:CFD10}
Let \(W\) be a finite-dimensional complex vector space, \(T\in
\operatorname{End}(W)\), and let \(W=\sum_iW_i\), with every
\(W_i\) stable under \(T\).  Then, for each \(\mu\in\mathbf C\),
\[
 W[\mu]_{\mathrm{gen}}
       =\sum_i\bigl(W_i\cap W[\mu]_{\mathrm{gen}}\bigr).
\]

If \(\mu\) is not an eigenvalue, both sides vanish.  Otherwise factor
the minimal polynomial as \(\prod_\nu(X-\nu)^{a_\nu}\).  The elementary
Chinese remainder theorem supplies \(p_\mu\in\mathbf C[X]\) that is one
modulo \((X-\mu)^{a_\mu}\) and zero modulo every other factor.
Consequently \(p_\mu(T)\) is the identity on the \(\mu\)-generalized
eigenspace and zero on the others.  Given a finite expression
\(w=\sum_iw_i\), apply \(p_\mu(T)\): stability keeps each projected
term in \(W_i\).  This proves the nontrivial inclusion; the other is
immediate.  No diagonalizability and no directness of \(\sum_iW_i\)
has been assumed.
% END CFD10

% BEGIN CFD11
\paragraph{CFD11. Refinement to the exact ordinary generalized eigenspaces.}
\label{claim:CFD11}
Assume \(e\ge c_0\), the original specialization assertions, and
\(C_{K,e}=\sum_i(Y_i\cap C_{K,e})\).  Then
\[
                         O_{K,e}=\sum_i(Y_i\cap O_{K,e}).
\]

The original \(O_{K,e}\) is the sum, inside \(C_{K,e}\), of the
generalized eigenspaces of \(T\) with
\(\lvert\iota^{-1}(\mu)\rvert_\ell=1\).  This identification is ASP6;
it involves no continuity of \(\iota\).  The original specialization
finiteness clause, proved in ASP10, makes \(C_{K,e}\) finite-dimensional.
CFD8 proves stability of both \(C_{K,e}\) and \(Y_i\cap C_{K,e}\).
Apply CFD10 to every eigenvalue with the indicated norm and add the
equalities.  Each resulting summand lies in \(Y_i\cap O_{K,e}\).
The reverse inclusion is automatic.  The argument includes nilpotent
Jordan parts and uses the source's unit condition, without replacing it
by a condition on the complex absolute value.
% END CFD11

% BEGIN CFD12
\paragraph{CFD12. Full original classical-parts endpoint and its predicate refinement.}
\label{claim:CFD12}
Under every hypothesis of the original \texttt{clParts\_alt}, for every
original \(K\) and \(e\ge c_0\),
\begin{align*}
 C_{K,e}&=\sum_{Y:\,P(\lambda,r_Y)}(Y\cap C_{K,e}),\\
 O_{K,e}&=\sum_{Y:\,P(\lambda,r_Y)}(Y\cap O_{K,e}).
\end{align*}
The same conclusion holds with a predicate on the whole part record,
provided the corresponding big-space decomposition is supplied.

For \(F\in C_{K,e}\), membership in \(H.\mathrm{Big}\) and the
given away level are among its defining equations.  CFD7 gives the
required Iwahori-fixedness.  The algebraic-representation assertion of
the original specialization datum supplies the exact algebraic-weight
hypothesis.  Apply \(H.\mathrm{AutomorphicHeckeSet}\) at this
\(\lambda,\xi,F\), obtaining membership in the stated sum of parts.
CFD9 preserves the original level and depth in that sum, proving the
first equality.  CFD11 proves the second.  The proof used the selecting
predicate only to specify the family being summed, so it proves the
predicate refinement as well.

This closes the original generic structural theorem, including every
level and the original map and depth assertions.  For the actual unitary
construction, FNL16 supplies \(H.\mathrm{Ax}\), and ASP13 supplies
the specialization datum; the automorphic part producer remains a
separate arithmetic input until it is explicitly constructed.  In
particular no nonvanishing of a selected residual Hecke block is inferred
from either equality when the block could be zero.
% END CFD12
